\documentclass[10pt,a4paper]{amsart}
\usepackage{amsmath,amssymb,amsthm,graphicx}
\usepackage{verbatim}
\newenvironment{reference}{\comment}{\endcomment}
%\newenvironment{reference}{}{}
\newenvironment{slogan}{\comment}{\endcomment}
\newenvironment{history}{\comment}{\endcomment}

% For commutative diagrams we use Xy-pic
\usepackage[all]{xy}

% We use 2cell for 2-commutative diagrams.
\xyoption{2cell}
% Theorem environments.
%
\theoremstyle{plain}
\newtheorem{theorem}[subsection]{Theorem}
\newtheorem{proposition}[subsection]{Proposition}
\newtheorem{lemma}[subsection]{Lemma}

\theoremstyle{definition}
\newtheorem{definition}[subsection]{Definition}
\newtheorem{example}[subsection]{Example}
\newtheorem{exercise}[subsection]{Exercise}
\newtheorem{situation}[subsection]{Situation}

\theoremstyle{remark}
\newtheorem{remark}[subsection]{Remark}
\newtheorem{remarks}[subsection]{Remarks}

\numberwithin{equation}{subsection}

% Macros
%
\def\lim{\mathop{\mathrm{lim}}\nolimits}
\def\colim{\mathop{\mathrm{colim}}\nolimits}
\def\Spec{\mathop{\mathrm{Spec}}}
\def\Hom{\mathop{\mathrm{Hom}}\nolimits}
\def\Ext{\mathop{\mathrm{Ext}}\nolimits}
\def\SheafHom{\mathop{\mathcal{H}\!\mathit{om}}\nolimits}
\def\SheafExt{\mathop{\mathcal{E}\!\mathit{xt}}\nolimits}
\def\Sch{\mathit{Sch}}
\def\Mor{\mathop{\mathrm{Mor}}\nolimits}
\def\Ob{\mathop{\mathrm{Ob}}\nolimits}
\def\Sh{\mathop{\mathit{Sh}}\nolimits}
\def\NL{\mathop{N\!L}\nolimits}
\def\CH{\mathop{\mathrm{CH}}\nolimits}
\def\proetale{{pro\text{-}\acute{e}tale}}
\def\etale{{\acute{e}tale}}
\def\QCoh{\mathit{QCoh}}
\def\Ker{\mathop{\mathrm{Ker}}}
\def\Im{\mathop{\mathrm{Im}}}
\def\Coker{\mathop{\mathrm{Coker}}}
\def\Coim{\mathop{\mathrm{Coim}}}

% Boxtimes
%
\DeclareMathSymbol{\boxtimes}{\mathbin}{AMSa}{"02}

%
% Macros for moduli stacks/spaces
%
\def\QCohstack{\mathcal{QC}\!\mathit{oh}}
\def\Cohstack{\mathcal{C}\!\mathit{oh}}
\def\Spacesstack{\mathcal{S}\!\mathit{paces}}
\def\Quotfunctor{\mathrm{Quot}}
\def\Hilbfunctor{\mathrm{Hilb}}
\def\Curvesstack{\mathcal{C}\!\mathit{urves}}
\def\Polarizedstack{\mathcal{P}\!\mathit{olarized}}
\def\Complexesstack{\mathcal{C}\!\mathit{omplexes}}
% \Pic is the operator that assigns to X its picard group, usage \Pic(X)
% \Picardstack_{X/B} denotes the Picard stack of X over B
% \Picardfunctor_{X/B} denotes the Picard functor of X over B
\def\Pic{\mathop{\mathrm{Pic}}\nolimits}
\def\Picardstack{\mathcal{P}\!\mathit{ic}}
\def\Picardfunctor{\mathrm{Pic}}
\def\Deformationcategory{\mathcal{D}\!\mathit{ef}}
\usepackage{cleveref}
\crefname{originalreference}{original reference}{original references}
\Crefname{originalreference}{Original reference}{Original references}
\makeatletter
\newlabel{situation-setup}{{02QL}{source}{}{sourcepointer-1}{}}
\newlabel{situation-setup@cref}{{[originalreference][0][]02QL}{[1][0][]source}{}{}{}}
\makeatother
\crefname{originallemma}{original lemma}{original lemmas}
\Crefname{originallemma}{Original lemma}{Original lemmas}
\makeatletter
\newlabel{more-algebra-lemma-base-change-lci-bis}{{0FV1}{source}{}{sourcepointer-2}{}}
\newlabel{more-algebra-lemma-base-change-lci-bis@cref}{{[originallemma][0][]0FV1}{[1][0][]source}{}{}{}}
\makeatother
\crefname{originallemma}{original lemma}{original lemmas}
\Crefname{originallemma}{Original lemma}{Original lemmas}
\makeatletter
\newlabel{more-morphisms-lemma-get-NL}{{0FV4}{source}{}{sourcepointer-3}{}}
\newlabel{more-morphisms-lemma-get-NL@cref}{{[originallemma][0][]0FV4}{[1][0][]source}{}{}{}}
\makeatother
\crefname{originallemma}{original lemma}{original lemmas}
\Crefname{originallemma}{Original lemma}{Original lemmas}
\makeatletter
\newlabel{morphisms-lemma-base-change-differentials}{{01V0}{source}{}{sourcepointer-4}{}}
\newlabel{morphisms-lemma-base-change-differentials@cref}{{[originallemma][0][]01V0}{[1][0][]source}{}{}{}}
\makeatother
\crefname{originallemma}{original lemma}{original lemmas}
\Crefname{originallemma}{Original lemma}{Original lemmas}
\makeatletter
\newlabel{morphisms-lemma-conormal-functorial-flat}{{0473}{source}{}{sourcepointer-5}{}}
\newlabel{morphisms-lemma-conormal-functorial-flat@cref}{{[originallemma][0][]0473}{[1][0][]source}{}{}{}}
\makeatother
\crefname{originallemma}{original lemma}{original lemmas}
\Crefname{originallemma}{Original lemma}{Original lemmas}
\makeatletter
\newlabel{lemma-gysin-excess}{{0FV7}{source}{}{sourcepointer-6}{}}
\newlabel{lemma-gysin-excess@cref}{{[originallemma][0][]0FV7}{[1][0][]source}{}{}{}}
\makeatother
\title{Base change of a Gysin map of a local complete intersection morphism is the top Chern class of the kernel of the conormal comparison}
\author{The Stacks Project Authors}
\date{stacks-project:chow.tex:local-snapshot-20261012; GNU FDL 1.2 or later}
\begin{document}
\maketitle
\noindent\textit{Editorial note.} All delivered mathematical text is the author's own, copied from the pinned Stacks Project chapter file chow.tex; no statement or proof step was rewritten, reordered or removed.\par\medskip
\noindent\textit{Editorial note.} Preamble provenance: the pinned chapter file opens with an input of the shared preamble and carries no preamble text of its own, so the project's single Stacks-wide preamble staged at sources/stacks-chow/preamble.tex is excerpted verbatim as two spans, lines 14-18 (Xy-pic setup) and lines 154-217 (the native theorem declarations and mathematical macros). Nothing was added, removed or redefined; the wrapper class and the amsmath/amssymb/amsthm/graphicx base packages are supplied by the assembler.\par\medskip
\noindent\textit{Editorial note.} The counted statement and proof contain no cross reference, no citation and no included file, so the delivered body is self contained.\par\medskip
\noindent\textit{Editorial note.} Canonical location of the counted proposition: Stacks tag 0FV9 (sources/stacks-chow/pinned-original-source.tex line 13738).\par\medskip
\noindent\textit{Allowed standard prerequisites (editorial).}\begin{itemize}
\item Finite local freeness of the kernel of the base changed conormal map. More on Algebra, Lemma labelled more-algebra-lemma-base-change-lci-bis, cited at proof lines 13763-13765. Provides the first assertion of the lemma.
\item Factorization of a local complete intersection morphism, the normal cone presentation and base change of differentials and conormal sheaves. More on Morphisms, Lemma labelled more-morphisms-lemma-get-NL, cited at proof lines 13778-13782; Morphisms, Lemmas labelled morphisms-lemma-base-change-differentials and morphisms-lemma-conormal-functorial-flat, cited at proof lines 13786-13792. Identifies the two descriptions of C and lets the excess computation be applied.
\item The excess formula for a regular immersion. Same chapter, Lemma labelled lemma-gysin-excess, cited at proof lines 13796-13798. Gives res(i\textasciicircum{}!) = c\_top(C\textasciicircum{}vee) o (i')\textasciicircum{}! for the immersion part.
\item Composition of the Gysin class of a local complete intersection morphism with the flat pullback. Same chapter, the identity f\textasciicircum{}! = i\textasciicircum{}! o h\textasciicircum{}* used at proof lines 13800-13805. Passes from the immersion to the full local complete intersection morphism.
\end{itemize}\medskip
\ifdefined\hypertarget\hypertarget{sourcepointer-1}{}\fi
\noindent\textit{Source reference.} Original reference 02QL: sources/stacks-chow/pinned-original-source.tex, source lines 1372--1372. Same-chapter pointer into the author environment situation at source line 1372; invoked at proof line None of the counted proof. It is retained only as a navigation pointer; every load-bearing step is inside the delivered span.\par\medskip
\ifdefined\hypertarget\hypertarget{sourcepointer-2}{}\fi
\noindent\textit{Source reference.} Original lemma 0FV1: sources/stacks-chow/tags.txt, source lines 19796--19796. Cross-chapter standard prerequisite located in the project tag list staged in the local snapshot; invoked at proof line 13765 of the counted proof. It is retained only as a navigation pointer; every load-bearing step is inside the delivered span.\par\medskip
\ifdefined\hypertarget\hypertarget{sourcepointer-3}{}\fi
\noindent\textit{Source reference.} Original lemma 0FV4: sources/stacks-chow/tags.txt, source lines 19799--19799. Cross-chapter standard prerequisite located in the project tag list staged in the local snapshot; invoked at proof line 13780 of the counted proof. It is retained only as a navigation pointer; every load-bearing step is inside the delivered span.\par\medskip
\ifdefined\hypertarget\hypertarget{sourcepointer-4}{}\fi
\noindent\textit{Source reference.} Original lemma 01V0: sources/stacks-chow/tags.txt, source lines 2429--2429. Cross-chapter standard prerequisite located in the project tag list staged in the local snapshot; invoked at proof line 13789 of the counted proof. It is retained only as a navigation pointer; every load-bearing step is inside the delivered span.\par\medskip
\ifdefined\hypertarget\hypertarget{sourcepointer-5}{}\fi
\noindent\textit{Source reference.} Original lemma 0473: sources/stacks-chow/tags.txt, source lines 5494--5494. Cross-chapter standard prerequisite located in the project tag list staged in the local snapshot; invoked at proof line 13792 of the counted proof. It is retained only as a navigation pointer; every load-bearing step is inside the delivered span.\par\medskip
\ifdefined\hypertarget\hypertarget{sourcepointer-6}{}\fi
\noindent\textit{Source reference.} Original lemma 0FV7: sources/stacks-chow/pinned-original-source.tex, source lines 12021--12021. Same-chapter pointer into the author environment lemma at source line 12021; invoked at proof line 13802 of the counted proof. It is retained only as a navigation pointer; every load-bearing step is inside the delivered span.\par\medskip
\begin{lemma}
\label{lemma-compare-gysin-base-change}
Let $(S, \delta)$ be as in Situation \ref{situation-setup}.
Consider a cartesian square
$$
\xymatrix{
X' \ar[d]_{f'} \ar[r]_{g'} &
X \ar[d]^f \\
Y' \ar[r]^g &
Y
}
$$
of schemes locally of finite type over $S$. Assume
\begin{enumerate}
\item both $f$ and $f'$ are local complete intersection morphisms, and
\item the Gysin map exists for $f$
\end{enumerate}
Then $\mathcal{C} = \Ker(H^{-1}((g')^*\NL_{X/Y}) \to H^{-1}(\NL_{X'/Y'}))$
is a finite locally free $\mathcal{O}_{X'}$-module, the Gysin map
exists for $f'$, and we have
$$
res(f^!) = c_{top}(\mathcal{C}^\vee) \circ (f')^!
$$
in $A^*(X' \to Y')$.
\end{lemma}

\begin{proof}
The fact that $\mathcal{C}$ is finite locally free follows immediately
from More on Algebra, Lemma \ref{more-algebra-lemma-base-change-lci-bis}.
Choose a factorization $f = h \circ i$ with $h : P \to Y$ smooth and $i$
an immersion. Then we can factor $f' = h' \circ i'$ where $h' : P' \to Y'$
and $i' : X' \to P'$ the base changes. Picture
$$
\xymatrix{
X' \ar[r] \ar[d] &
P' \ar[r] \ar[d] &
Y' \ar[d] \\
X \ar[r]^i &
P \ar[r]^h &
Y
}
$$
In particular, we see that the Gysin map exists for $f'$. By
More on Morphisms, Lemmas \ref{more-morphisms-lemma-get-NL}
we have
$$
\NL_{X/Y} = \left( \mathcal{C}_{X/P} \to i^*\Omega_{P/Y} \right)
$$
where $\mathcal{C}_{X/P}$ is the conormal sheaf of the embedding $i$.
Similarly for the primed version. We have
$(g')^*i^*\Omega_{P/Y} = (i')^*\Omega_{P'/Y'}$ because
$\Omega_{P/Y}$ pulls back to $\Omega_{P'/Y'}$ by
Morphisms, Lemma \ref{morphisms-lemma-base-change-differentials}.
Also, recall that $(g')^*\mathcal{C}_{X/P} \to \mathcal{C}_{X'/P'}$
is surjective, see
Morphisms, Lemma \ref{morphisms-lemma-conormal-functorial-flat}.
We deduce that the sheaf $\mathcal{C}$ is canonicallly
isomorphic to the kernel of the map
$(g')^*\mathcal{C}_{X/P} \to \mathcal{C}_{X'/P'}$
of finite locally free modules. Recall that $i^!$ is defined
using $\mathcal{N} = \mathcal{C}_{Z/X}^\vee$ and similarly
for $(i')^!$. Thus we have
$$
res(i^!) = c_{top}(\mathcal{C}^\vee) \circ (i')^!
$$
in $A^*(X' \to P')$ by an application of Lemma \ref{lemma-gysin-excess}.
Since finally we have $f^! = i^! \circ h^*$,
$(f')^! = (i')^! \circ (h')^*$, and $(h')^* = res(h^*)$ we conclude.
\end{proof}

\end{document}
